% from main.tex
%%%%%%%%%%  scan idx 237  |  folio 224  %%%%%%%%%%
\textbf{Remarque 1.1.2.} Le lemme précédent s'applique notamment au cas où \(i\) est une immersion régulière (VII~1.4). En effet, on a localement, grâce au complexe de Kozsul, une résolution gauche finie de \(i_*(\mathcal O_X)\) par des \(\mathcal O_Y\)-Modules libres de type fini, donc \(i_*(\mathcal O_X)\) est parfait (I 4.8 et 4.14) (donc a fortiori pseudo-cohérent).

\textbf{Lemme 1.1.3.} Soit un carré cartésien de schémas
\[
\begin{tikzcd}
X \arrow[d,"i"'] & X' \arrow[l,"f"'] \arrow[d,"i'"]\\
Y & Y'{,} \arrow[l,"g"']
\end{tikzcd}
\]
où \(i\) est une immersion fermée, et soit \(E\) un complexe de \(\mathcal O_X\)-Modules. Alors la flèche canonique de changement de base (EGA~III 1.4.15.4)
\[
  g^*i_*(E)\longrightarrow i'_*f^*(E)
\]
est un isomorphisme.

(ii) Posons \(E'=f^*(E)\). Si \(g\) est plat et \(i_*(E)\) \(n\)-pseudo-cohérent, alors \(i'_*(E')\) est \(n\)-pseudo-cohérent. Inversement, si \(g\) est fidèlement plat et \(i'_*(E')\) \(n\)-pseudo-cohérent, alors \(i_*(E)\) est \(n\)-pseudo-cohérent.

\emph{Preuve.} L'assertion se vérifie trivialement sur les fibres.

\textbf{Lemme 1.1.4.}
